%! Author = zhiweizhang
%! Date = 10.06.25

\chapter{Ethical Considerations}\label{chapter:ethical-considerations}

The development of \ac{ML} systems for the classification of Nazi symbols entails a range of ethical
challenges due to the sensitive, harmful, and potentially illegal nature of the target content. These concerns
reflect broader societal expectations for responsible AI, as articulated in international frameworks such as \ac{UNESCO}’s
\textit{Recommendation on the Ethics of Artificial Intelligence}~\parencite{unesco2021}. This chapter
outlines the core ethical principles guiding the research, grounded in widely recognized frameworks such as the \ac{IEEE}
\textit{Ethically Aligned Design} guidelines~\parencite{ieee2021ethics}, the EU AI Act~\parencite{euai2023}, and the
\ac{DFG}’s \textit{Guidelines for Safeguarding Good Research Practice}~\parencite{dfg2022}. These considerations are
essential for ensuring that the study is conducted responsibly, complies with legal norms, and anticipates the
social implications of deploying such a system.

\begin{itemize}
    \item \textbf{Foundational Ethical Frameworks:} This research is informed by established ethical guidelines such
    as the IEEE Ethically Aligned Design~\parencite{ieee2021ethics} and the Belmont Report~\parencite{belmont1979},
    the latter—originally developed for biomedical research—offers enduring principles like
    respect for persons and beneficence that extend to socially impactful AI systems. These principles guide our
    handling of sensitive data and system design choices.

    \item \textbf{Handling of Sensitive Content:} Nazi symbols are strongly associated with hate, violence, and
    genocide. Datasets containing such imagery must be handled with care to avoid trivializing or normalizing hateful
    ideologies. Access to these datasets is strictly limited, and appropriate content warnings and disclaimers are used.
    Data management follows established protocols for working with sensitive or controversial materials, in accordance
    with research integrity standards outlined by the \ac{DFG}~\parencite{dfg2022}.

    \item \textbf{Purpose and Intent:} The goal of this research is to assist in identifying and moderating harmful
    content, not to promote or propagate extremist ideologies. Transparency in the project's objectives helps
    prevent misuse, aligning with principles of socially beneficial AI outlined by the \ac{IEEE} and
    \ac{UNESCO}~\parencite{ieee2021ethics,unesco2021}.

    \item \textbf{Legal Compliance:} In countries such as Germany and Austria, strict laws (e.g., Section 86a of the
    German Criminal Code) prohibit the display and distribution of Nazi symbols. This study complies with all
    applicable regulations, ensuring that content is accessed solely for legitimate research purposes and not
    redistributed~\parencite{bundesregierung2023}.

    \item \textbf{Data Sourcing and Annotation:} All datasets were collected from publicly available
    repositories—primarily Kaggle and Roboflow—which permit academic and non-commercial use under their respective
    licenses\footnote{Kaggle Terms of Service: \url{https://www.kaggle.com/terms}; Roboflow Terms:
    \url{https://roboflow.com/terms}}. No restricted or proprietary content was used. All annotations were conducted
    manually by the researcher. Content warnings were observed throughout the process to mitigate potential
    psychological distress, as also described in ~\autofullref{sec:data-construction} from ~\autofullref{chapter:methodology}.

    \item \textbf{Bias and Misclassification Risks:} Classification systems may mistakenly flag benign content or
    overlook altered forms of hate symbols. Such errors could lead to unjust moderation or missed threats. To
    address this, the study prioritized balanced training data, fairness-aware evaluation, and extensive model
    validation in line with AI fairness principles~\parencite{mehrabi2021survey}.

    \item \textbf{Dual-Use Concerns:} While developed for moderation purposes, such a classification system could be
    repurposed for malicious use, such as evading content filters or enabling authoritarian surveillance.
    Recognizing this dual-use risk~\parencite{brundage2018malicious}, the research limits model distribution and
    recommends transparency, access control, and red-teaming before deployment.

    \item \textbf{Privacy and Consent:} If any images included identifiable individuals, anonymization procedures
    were applied in accordance with \ac{GDPR} principles and ethical AI guidelines~\parencite{gdpr2016}. No \ac{PII}
    was intentionally collected or retained.

    \item \textbf{Transparency and Accountability:} All stages of the research—including data provenance,
    preprocessing, model selection, and limitations—are documented and disclosed. This transparency aligns with
    reproducible research norms and promotes accountability in the deployment of AI systems on sensitive
    content~\parencite{Mitchell_2019}.
\end{itemize}

\noindent
In sum, these ethical frameworks are not peripheral to the study but foundational to its design and execution. They
inform not only how data is sourced, labeled, and secured, but also how models are evaluated, interpreted, and
constrained in their use. With these considerations in place, the following chapter details the methodological
pipeline through which the classification system is developed and validated.
